\documentclass[aspectratio=169]{beamer}
\usetheme{Madrid}
\usecolortheme{default}
\usepackage[utf8]{inputenc}
\usepackage{amsmath}
\usepackage{amsfonts}
\usepackage{amssymb}
\usepackage{enumerate}

\setbeamertemplate{navigation symbols}{}
\setbeamerfont{block body}{size=\small}
\setbeamerfont{alertblock body}{size=\small}

\title{CRISC Mixed Practice Exam II -- 100 Real-Life Questions}
\subtitle{Weighted to the Four CRISC Domains -- Practical Scenarios}
\author{CRISC Mastery Series}
\date{\today}

% ------------------------------------------------------------
% \hq{number}{topic}{question}{options}{letter}{explanation}{takeaway}
% Creates one question frame and one explanation frame.
% ------------------------------------------------------------
\newcommand{\hq}[7]{%
	\begin{frame}{Q#1: #2}
		\begin{block}{Question}
			#3
		\end{block}
		\begin{itemize}
			#4
		\end{itemize}
	\end{frame}
	\begin{frame}{Q#1: Explanation}
		\begin{block}{Correct Answer: #5}
			\begin{itemize}
				#6
			\end{itemize}
		\end{block}
		\begin{alertblock}{Key Takeaway}
			#7
		\end{alertblock}
	\end{frame}
}

\begin{document}

	\begin{frame}
		\titlepage
	\end{frame}

	\begin{frame}{Instructions}
		\begin{alertblock}{How to Use This Deck}
			Think like a risk practitioner in a real organization.\\
			Focus on business context, ownership, and practical decision-making.\\
			The answer and a key takeaway are on the slide following each question.
		\end{alertblock}
	\end{frame}

	\begin{frame}{Outline}
		\tableofcontents
	\end{frame}

	% ============================================================
	\section{Governance (26 percent of the exam)}
	% ============================================================

	\hq{1}{Risk Policy Contents}
	{Which element is \textbf{MOST} essential in an enterprise risk management policy?}
	{\item[\textbf{(A)}] A list of approved software
		\item[\textbf{(B)}] Employee dress standards
		\item[\textbf{(C)}] Office seating plans
		\item[\textbf{(D)}] Defined roles, risk appetite alignment, and required processes for identifying, assessing, responding to, and reporting risk}
	{D}
	{\item Policy sets the framework for risk activities.
		\item Roles create accountability.
		\item Alignment to appetite guides decisions.}
	{A risk policy defines who does what, and to what standard.}

	\hq{2}{CISO Reporting Line}
	{The CISO reports to the CIO, and security concerns are often overridden for cost reasons. What is the \textbf{BEST} governance improvement?}
	{\item[\textbf{(A)}] Remove the CISO role
		\item[\textbf{(B)}] Make the CISO report to the help desk
		\item[\textbf{(C)}] Hide security concerns
		\item[\textbf{(D)}] Give the CISO a reporting line or access that supports independence, such as direct access to senior executives or the board}
	{D}
	{\item Independence reduces conflicts of interest.
		\item Senior access ensures material risks are heard.
		\item Trade-offs should be visible to decision makers.}
	{Risk voices need a path to the top.}

	\hq{3}{Key Person Dependency}
	{Only one engineer understands a critical system. What is the \textbf{BEST} risk response?}
	{\item[\textbf{(A)}] Accept the risk forever with no action
		\item[\textbf{(B)}] Document the system, cross-train staff, and establish succession and backup coverage
		\item[\textbf{(C)}] Prevent the engineer from taking vacation
		\item[\textbf{(D)}] Delete the system documentation}
	{B}
	{\item Single points of failure include people.
		\item Knowledge sharing reduces dependency.
		\item Succession planning supports resilience.}
	{People can be single points of failure too.}

	\hq{4}{Tone at the Top}
	{Executives routinely bypass security controls for their own convenience. What is the \textbf{MOST} significant effect?}
	{\item[\textbf{(A)}] It strengthens controls
		\item[\textbf{(B)}] It undermines the risk culture because staff follow leaders' behavior rather than written policy
		\item[\textbf{(C)}] It has no effect
		\item[\textbf{(D)}] It improves audit results}
	{B}
	{\item Leaders shape norms.
		\item Exemptions without approval erode trust.
		\item Consistent standards support compliance.}
	{Staff copy what leaders do.}

	\hq{5}{Management Override}
	{Which control is \textbf{MOST} effective against management override of financial and IT controls?}
	{\item[\textbf{(A)}] Trusting executives completely
		\item[\textbf{(B)}] Eliminating all controls
		\item[\textbf{(C)}] Independent review, monitoring of overrides, and a whistleblower channel
		\item[\textbf{(D)}] Annual training only}
	{C}
	{\item Override is a known fraud risk.
		\item Independent oversight detects abuse.
		\item Reporting channels expose violations.}
	{Watch the people who can bypass the controls.}

	\hq{6}{Whistleblower Policy}
	{What is the primary purpose of a whistleblower policy?}
	{\item[\textbf{(A)}] To enable staff to report misconduct or risk concerns safely, without retaliation
		\item[\textbf{(B)}] To punish complainers
		\item[\textbf{(C)}] To reduce paperwork
		\item[\textbf{(D)}] To publicize employee names}
	{A}
	{\item Safe reporting surfaces hidden risk.
		\item Anti-retaliation protection builds trust.
		\item Independent handling ensures fairness.}
	{Safe reporting reveals risk.}

	\hq{7}{ISO 31000}
	{An organization wants a general, industry-neutral guideline for risk management principles and process. Which standard \textbf{BEST} fits?}
	{\item[\textbf{(A)}] A cabling specification
		\item[\textbf{(B)}] A coding language standard
		\item[\textbf{(C)}] A paper size standard
		\item[\textbf{(D)}] ISO 31000}
	{D}
	{\item ISO 31000 provides principles and a process.
		\item It applies to many organization types.
		\item It complements, not replaces, domain frameworks.}
	{ISO 31000 is the general risk management guideline.}

	\hq{8}{ISO 27001}
	{Management wants an internationally recognized, certifiable approach for establishing an information security management system. Which standard fits?}
	{\item[\textbf{(A)}] ISO 9660
		\item[\textbf{(B)}] ISO 8601
		\item[\textbf{(C)}] ISO 4217
		\item[\textbf{(D)}] ISO/IEC 27001}
	{D}
	{\item 27001 specifies ISMS requirements.
		\item Certification is by independent bodies.
		\item It is risk-based.}
	{27001 is the certifiable ISMS standard.}

	\hq{9}{COSO Concept}
	{Which framework is \textbf{MOST} associated with enterprise risk management integrated with strategy and performance?}
	{\item[\textbf{(A)}] A network protocol standard
		\item[\textbf{(B)}] A programming framework
		\item[\textbf{(C)}] COSO Enterprise Risk Management framework
		\item[\textbf{(D)}] A building code}
	{C}
	{\item COSO ERM links risk to strategy.
		\item It is widely used for governance.
		\item It complements IT frameworks.}
	{COSO ERM ties risk to strategy and performance.}

	\hq{10}{Quantitative Framework}
	{Which framework is specifically designed for quantifying information risk in financial terms?}
	{\item[\textbf{(A)}] A color-coding scheme
		\item[\textbf{(B)}] A password standard
		\item[\textbf{(C)}] Factor Analysis of Information Risk (FAIR)
		\item[\textbf{(D)}] A seating chart}
	{C}
	{\item FAIR models loss event frequency and magnitude.
		\item Outputs support cost-benefit analysis.
		\item It needs sound input data.}
	{FAIR quantifies risk in dollars.}

	\hq{11}{Maturity Targets}
	{Management sets a maturity target of the highest level for every process. What is the \textbf{BEST} challenge?}
	{\item[\textbf{(A)}] Agree to the highest level always
		\item[\textbf{(B)}] Set targets based on business need and risk, since the highest maturity everywhere may not be cost-justified
		\item[\textbf{(C)}] Remove all targets
		\item[\textbf{(D)}] Use random targets}
	{B}
	{\item Maturity should match value and risk.
		\item Resources are limited.
		\item Priorities should be explicit.}
	{Target maturity by need, not by ambition.}

	\hq{12}{Benchmarking}
	{A board compares its security spending to a peer average. What is the \textbf{PRIMARY} limitation?}
	{\item[\textbf{(A)}] Benchmarks are always illegal
		\item[\textbf{(B)}] Peers have different risk profiles, so benchmarks do not show whether spending matches this organization's risk
		\item[\textbf{(C)}] Peers never spend anything
		\item[\textbf{(D)}] Benchmarks are perfectly precise}
	{B}
	{\item Risk appetite and exposure differ.
		\item Benchmarks are useful context only.
		\item Decisions should rest on own risk.}
	{Benchmarks inform; risk decides.}

	\hq{13}{Policy Sunset}
	{A policy has no owner or review date. What is the \textbf{BEST} governance action?}
	{\item[\textbf{(A)}] Ignore it
		\item[\textbf{(B)}] Delete it without review
		\item[\textbf{(C)}] Assign an owner and set a review cycle, then retire or update the policy accordingly
		\item[\textbf{(D)}] Enforce it as is}
	{C}
	{\item Ownership drives maintenance.
		\item Review dates keep policies current.
		\item Retirement avoids conflicting guidance.}
	{Every policy needs an owner and a review date.}

	\hq{14}{Regulatory Examination Readiness}
	{A regulatory examination is scheduled in three months. What should the risk practitioner do?}
	{\item[\textbf{(A)}] Assess readiness, close known gaps, and ensure evidence is organized and accurate
		\item[\textbf{(B)}] Wait for the examiners to arrive
		\item[\textbf{(C)}] Alter documents to look better
		\item[\textbf{(D)}] Avoid preparation}
	{A}
	{\item Preparation reduces findings.
		\item Evidence must be truthful.
		\item Known gaps should be disclosed with plans.}
	{Be ready, be honest.}

	\hq{15}{Records Retention Governance}
	{Departments keep records under inconsistent retention rules. What should governance implement?}
	{\item[\textbf{(A)}] A records retention schedule based on legal and business requirements, with defined disposal procedures
		\item[\textbf{(B)}] Keep everything forever
		\item[\textbf{(C)}] Delete records daily
		\item[\textbf{(D)}] Let each person decide}
	{A}
	{\item Retention reduces legal and storage risk.
		\item Legal obligations differ by record type.
		\item Disposal should be secure.}
	{Retain what is required, dispose of the rest.}

	\hq{16}{Cyber Expertise on the Board}
	{The board lacks cyber risk knowledge. What is the \textbf{BEST} approach?}
	{\item[\textbf{(A)}] Provide education and use qualified advisors or committees to support informed oversight
		\item[\textbf{(B)}] Ignore cyber topics
		\item[\textbf{(C)}] Delegate all decisions to IT
		\item[\textbf{(D)}] Avoid reporting cyber risk}
	{A}
	{\item Oversight requires understanding.
		\item Advisors provide expertise.
		\item Education builds long-term capability.}
	{Boards must understand enough to oversee.}

	\hq{17}{Balanced Scorecard}
	{How does a balanced scorecard help IT governance?}
	{\item[\textbf{(A)}] It lists server names
		\item[\textbf{(B)}] It replaces strategy
		\item[\textbf{(C)}] It prevents attacks
		\item[\textbf{(D)}] It links IT performance measures to business objectives across financial, customer, process, and learning perspectives}
	{D}
	{\item It translates strategy into measures.
		\item It balances short- and long-term views.
		\item It supports alignment.}
	{Scorecards align IT measures to strategy.}

	\hq{18}{Risk Acceptance Register}
	{Why maintain a central register of accepted risks?}
	{\item[\textbf{(A)}] To provide transparency, ensure periodic review, and keep accountability for accepted exposures
		\item[\textbf{(B)}] To hide risks from auditors
		\item[\textbf{(C)}] To avoid reviews
		\item[\textbf{(D)}] To remove owners}
	{A}
	{\item Accepted risks should not be forgotten.
		\item Reviews confirm continued validity.
		\item Visibility supports governance.}
	{Accepted does not mean forgotten.}

	\hq{19}{Anti-Fraud Controls}
	{Which IT control \textbf{BEST} supports fraud prevention in financial systems?}
	{\item[\textbf{(A)}] Segregation of duties, access reviews, and transaction monitoring
		\item[\textbf{(B)}] Shared administrator passwords
		\item[\textbf{(C)}] Unlimited access
		\item[\textbf{(D)}] No logging}
	{A}
	{\item Fraud exploits access and weak oversight.
		\item Monitoring detects anomalies.
		\item Separation limits abuse.}
	{Limit access, separate duties, watch transactions.}

	\hq{20}{Third-Party Ethics}
	{A key supplier is accused of using forced labor. What is the \textbf{BEST} first step?}
	{\item[\textbf{(A)}] Ignore the allegation
		\item[\textbf{(B)}] Publicly condemn without facts
		\item[\textbf{(C)}] Assess the credibility and impact, engage the vendor, and escalate through governance and compliance channels
		\item[\textbf{(D)}] Hide the information}
	{C}
	{\item Ethical issues create reputational and legal risk.
		\item Fact-finding comes before decisions.
		\item Governance decides the response.}
	{Investigate, then escalate.}

	\hq{21}{Governance Metrics}
	{Which metric \textbf{BEST} indicates effective risk governance?}
	{\item[\textbf{(A)}] Number of pages in reports
		\item[\textbf{(B)}] Number of committee meetings
		\item[\textbf{(C)}] Timely risk decisions made by accountable owners within appetite, and acted-on escalations
		\item[\textbf{(D)}] Number of risks hidden}
	{C}
	{\item Outcomes matter more than activity.
		\item Accountability shows maturity.
		\item Escalations demonstrate functioning oversight.}
	{Governance is judged by decisions and follow-through.}

	\hq{22}{Strategy and Technology Debt}
	{Years of postponed upgrades have created large technology debt. How should this be governed?}
	{\item[\textbf{(A)}] Ignore it
		\item[\textbf{(B)}] Hide it from the board
		\item[\textbf{(C)}] Treat it as an IT-only problem
		\item[\textbf{(D)}] Track technical debt as a risk with owners and funding decisions aligned to appetite}
	{D}
	{\item Debt increases failure and security risk.
		\item Owners and plans are needed.
		\item Funding is a business decision.}
	{Technical debt is business risk.}

	\hq{23}{Risk Practitioner Independence}
	{A risk practitioner is evaluated by the manager of the department being assessed. What is the concern?}
	{\item[\textbf{(A)}] None, all is fine
		\item[\textbf{(B)}] Practitioners should avoid assessments
		\item[\textbf{(C)}] Managers cannot evaluate anyone
		\item[\textbf{(D)}] Objectivity may be compromised, so the reporting line should support independence}
	{D}
	{\item Evaluations can pressure ratings.
		\item Independent reporting protects objectivity.
		\item Governance should address it.}
	{Do not assess the person who controls your career.}

	\hq{24}{Time-Bound Governance Decisions}
	{Governance decisions on risk responses take months, leaving risks exposed. What is the \textbf{BEST} fix?}
	{\item[\textbf{(A)}] Wait longer
		\item[\textbf{(B)}] Skip decisions
		\item[\textbf{(C)}] Define decision timelines and delegated authority based on risk level
		\item[\textbf{(D)}] Ban governance meetings}
	{C}
	{\item Timeliness affects exposure.
		\item Delegation speeds routine decisions.
		\item Higher risks still escalate.}
	{Governance must move at the speed of risk.}

	\hq{25}{Risk Appetite Statement Quality}
	{Which risk appetite statement is \textbf{BEST} for operational use?}
	{\item[\textbf{(A)}] We want to be safe
		\item[\textbf{(B)}] We tolerate no more than four hours of outage for customer-facing systems and zero unmitigated critical vulnerabilities older than 30 days
		\item[\textbf{(C)}] We avoid all risk
		\item[\textbf{(D)}] We will do our best}
	{B}
	{\item Good statements are specific and measurable.
		\item They can be monitored with indicators.
		\item Vague statements cannot guide decisions.}
	{Appetite statements must be measurable.}

	\hq{26}{Linking Risk to Performance}
	{How can risk management \textbf{BEST} be integrated with business performance management?}
	{\item[\textbf{(A)}] By keeping risk separate
		\item[\textbf{(B)}] By including risk indicators and appetite in planning, performance reviews, and decision making
		\item[\textbf{(C)}] By reporting risk only to audit
		\item[\textbf{(D)}] By ignoring risk during planning}
	{B}
	{\item Integrated risk supports strategy.
		\item Indicators inform trade-offs.
		\item Decisions should reflect appetite.}
	{Risk belongs in planning and performance conversations.}

	% ============================================================
	\section{Risk Assessment (20 percent of the exam)}
	% ============================================================

	\hq{27}{Delphi Technique}
	{Which technique collects anonymous expert estimates over several rounds to reach consensus?}
	{\item[\textbf{(A)}] The Delphi technique
		\item[\textbf{(B)}] A brainstorm among executives only
		\item[\textbf{(C)}] A random poll
		\item[\textbf{(D)}] A vendor presentation}
	{A}
	{\item Anonymity reduces dominance bias.
		\item Iteration refines estimates.
		\item Consensus improves reliability.}
	{Delphi means anonymous rounds of expert input.}

	\hq{28}{Bow-Tie Analysis}
	{What does a bow-tie diagram help the risk practitioner visualize?}
	{\item[\textbf{(A)}] Causes, the central risk event, consequences, and the preventive and recovery controls between them
		\item[\textbf{(B)}] Employee schedules
		\item[\textbf{(C)}] Network cable routes
		\item[\textbf{(D)}] Financial statements only}
	{A}
	{\item It links causes to consequences.
		\item Preventive controls sit on the cause side.
		\item Recovery controls sit on the consequence side.}
	{Bow-ties show causes, event, consequences, and controls.}

	\hq{29}{Attack Frameworks}
	{A team wants a common catalog of adversary tactics and techniques to design detections and assess coverage. What \textbf{BEST} fits?}
	{\item[\textbf{(A)}] A purchasing catalog
		\item[\textbf{(B)}] A chart of accounts
		\item[\textbf{(C)}] A coding style guide
		\item[\textbf{(D)}] The MITRE ATT\&CK knowledge base}
	{D}
	{\item It structures real-world attacker behavior.
		\item Coverage mapping highlights gaps.
		\item It supports threat-informed defense.}
	{ATT\&CK maps adversary behavior to detections.}

	\hq{30}{Kill Chain}
	{Which concept describes the stages of an attack from reconnaissance to actions on objectives, helping defenders find opportunities to disrupt?}
	{\item[\textbf{(A)}] The cyber kill chain
		\item[\textbf{(B)}] A supply chain schedule
		\item[\textbf{(C)}] A reporting chain of command
		\item[\textbf{(D)}] A key chain inventory}
	{A}
	{\item Stages show where to detect or block.
		\item Earlier disruption reduces impact.
		\item Defense in depth maps to stages.}
	{Break the chain early.}

	\hq{31}{Likelihood Scale Definition}
	{Assessors interpret terms like unlikely and possible differently. What is the \textbf{BEST} improvement?}
	{\item[\textbf{(A)}] Use vague words only
		\item[\textbf{(B)}] Let each assessor decide
		\item[\textbf{(C)}] Use only the highest level
		\item[\textbf{(D)}] Define each likelihood level with clear descriptions, such as ranges for frequency or probability}
	{D}
	{\item Clear definitions reduce subjectivity.
		\item Consistency enables comparison.
		\item Scales should be approved.}
	{Define the scale, then use it.}

	\hq{32}{Control Effectiveness Rating}
	{How should control effectiveness be determined for inclusion in a risk assessment?}
	{\item[\textbf{(A)}] By evidence from testing, monitoring, and incidents rather than by design intent alone
		\item[\textbf{(B)}] By the vendor's brochure
		\item[\textbf{(C)}] By the control owner's confidence
		\item[\textbf{(D)}] By the control's price}
	{A}
	{\item Evidence shows real performance.
		\item Design and operation both matter.
		\item Re-evaluate as conditions change.}
	{Rate controls on evidence.}

	\hq{33}{Geopolitical Risk}
	{Cloud services are hosted in a region facing increasing political instability. What should the assessment consider?}
	{\item[\textbf{(A)}] Impact on availability, legal access to data, and the feasibility of relocating services
		\item[\textbf{(B)}] Nothing; politics is not a risk
		\item[\textbf{(C)}] Only the hosting price
		\item[\textbf{(D)}] Only the vendor's logo}
	{A}
	{\item External events can disrupt services.
		\item Legal access and sanctions may matter.
		\item Exit options reduce exposure.}
	{External events belong in the assessment.}

	\hq{34}{Technical Debt Assessment}
	{How should the risk from aging, unsupported systems be assessed?}
	{\item[\textbf{(A)}] Consider age only
		\item[\textbf{(B)}] Ignore business criticality
		\item[\textbf{(C)}] Consider only the color of the casing
		\item[\textbf{(D)}] Consider vulnerability exposure, business criticality, supportability, and the cost and time to replace}
	{D}
	{\item Multiple factors determine risk.
		\item Criticality sets priority.
		\item Replacement timelines inform response.}
	{Age alone is not risk; exposure and criticality are.}

	\hq{35}{Data Inventory}
	{Why is a data inventory or data map important for risk assessment?}
	{\item[\textbf{(A)}] It lists office plants
		\item[\textbf{(B)}] It replaces classification
		\item[\textbf{(C)}] It shows where sensitive data resides and flows, enabling accurate impact and exposure assessment
		\item[\textbf{(D)}] It eliminates privacy law}
	{C}
	{\item You cannot protect unknown data.
		\item Flows reveal third-party exposure.
		\item It supports privacy obligations.}
	{Know where the data is.}

	\hq{36}{Phishing Risk Assessment}
	{Which approach \textbf{BEST} assesses the organization's phishing risk?}
	{\item[\textbf{(A)}] Count the number of employees
		\item[\textbf{(B)}] Combine simulated phishing results, incident history, email controls, and training effectiveness
		\item[\textbf{(C)}] Check office wall colors
		\item[\textbf{(D)}] Review only the firewall}
	{B}
	{\item Multiple indicators give a fuller picture.
		\item Simulations measure behavior.
		\item Controls reduce impact.}
	{Use several lenses on human risk.}

	\hq{37}{M\&A Due Diligence}
	{Before acquiring a company, which risk activity is \textbf{MOST} appropriate?}
	{\item[\textbf{(A)}] Skip it to speed the deal
		\item[\textbf{(B)}] Review only the logo
		\item[\textbf{(C)}] Ask the target to self-certify
		\item[\textbf{(D)}] Perform cyber and IT due diligence on the target's security posture, incidents, and liabilities}
	{D}
	{\item Acquired weaknesses become your risk.
		\item Findings may affect valuation or terms.
		\item Integration planning depends on results.}
	{You inherit what you acquire.}

	\hq{38}{Fault Tree Analysis}
	{What does fault tree analysis do?}
	{\item[\textbf{(A)}] Predicts stock prices
		\item[\textbf{(B)}] Schedules meetings
		\item[\textbf{(C)}] Counts customer complaints
		\item[\textbf{(D)}] Works backward from an undesired event to identify the combinations of failures that could cause it}
	{D}
	{\item It decomposes causes logically.
		\item It reveals single points of failure.
		\item It supports design improvement.}
	{Fault trees trace causes back from the failure.}

	\hq{39}{Risk Correlation}
	{Two supposedly independent risks always materialize together. What does this suggest?}
	{\item[\textbf{(A)}] It is a coincidence
		\item[\textbf{(B)}] They are correlated through a common cause, so combined impact should be assessed
		\item[\textbf{(C)}] The risks are unrelated
		\item[\textbf{(D)}] The risks can be removed}
	{B}
	{\item Correlation increases aggregate exposure.
		\item Common causes need shared responses.
		\item Independent assumptions may be wrong.}
	{Correlated risks hit together.}

	\hq{40}{Outsourced SOC Assessment}
	{A managed security operations provider monitors your environment. What should the assessment confirm?}
	{\item[\textbf{(A)}] Only the logo of the provider
		\item[\textbf{(B)}] Only the price
		\item[\textbf{(C)}] Nothing; outsourcing removes risk
		\item[\textbf{(D)}] Coverage, detection quality, escalation and response times, and the provider's own control assurance}
	{D}
	{\item You depend on their effectiveness.
		\item Response times affect impact.
		\item Independent assurance provides evidence.}
	{Outsourced monitoring still needs oversight.}

	\hq{41}{Insider Risk Assessment}
	{Which factor is \textbf{MOST} relevant when assessing insider risk?}
	{\item[\textbf{(A)}] The employee's favorite color
		\item[\textbf{(B)}] Privileged access, sensitive data exposure, monitoring coverage, and organizational stressors
		\item[\textbf{(C)}] Office size
		\item[\textbf{(D)}] Weather}
	{B}
	{\item Access and data determine potential impact.
		\item Monitoring affects detection.
		\item Stressors can influence behavior.}
	{Insider risk is access plus opportunity plus motive.}

	\hq{42}{Risk Appetite in Scenarios}
	{How should risk appetite be used when evaluating a scenario's result?}
	{\item[\textbf{(A)}] Ignore appetite
		\item[\textbf{(B)}] Change appetite to match the result
		\item[\textbf{(C)}] Compare the assessed risk to appetite and tolerance to determine whether response is needed
		\item[\textbf{(D)}] Use it only after incidents}
	{C}
	{\item Appetite is the benchmark for acceptability.
		\item Tolerance gives measurable thresholds.
		\item Response decisions follow.}
	{Assess against appetite.}

	\hq{43}{RTO Conflicts}
	{A critical application has an RTO of 2 hours, but its database dependency has a 12-hour recovery. What should the assessment conclude?}
	{\item[\textbf{(A)}] The RTO is satisfied
		\item[\textbf{(B)}] The recovery objectives are inconsistent, so the dependency must be improved or the RTO revisited
		\item[\textbf{(C)}] The dependency does not matter
		\item[\textbf{(D)}] Both can be ignored}
	{B}
	{\item The chain is only as fast as its slowest dependency.
		\item Gaps need response.
		\item Business must understand the real RTO.}
	{Recovery is limited by the slowest dependency.}

	\hq{44}{Control Dependency Failure}
	{Several controls rely on the same authentication service. What risk does this reveal?}
	{\item[\textbf{(A)}] A common-mode failure, where one failure could disable multiple controls
		\item[\textbf{(B)}] No risk
		\item[\textbf{(C)}] Better redundancy
		\item[\textbf{(D)}] Faster recovery}
	{A}
	{\item Shared dependencies reduce independence.
		\item Defense in depth needs diversity.
		\item Resilience of the shared service matters.}
	{Layers must not share one point of failure.}

	\hq{45}{Scoping an Enterprise Assessment}
	{A new enterprise risk assessment is requested. What should be defined \textbf{FIRST}?}
	{\item[\textbf{(A)}] The font of the report
		\item[\textbf{(B)}] Scope, objectives, criteria, and stakeholders, aligned with business priorities
		\item[\textbf{(C)}] The final conclusions
		\item[\textbf{(D)}] Nothing; begin immediately}
	{B}
	{\item Scope determines effort and relevance.
		\item Criteria ensure consistency.
		\item Stakeholders provide input and ownership.}
	{Scope before scoring.}

	\hq{46}{Assessment Documentation}
	{A risk assessment is complete. Which documentation practice is \textbf{BEST}?}
	{\item[\textbf{(A)}] Keep conclusions only
		\item[\textbf{(B)}] Delete the working papers
		\item[\textbf{(C)}] Record scope, methodology, assumptions, data sources, results, and limitations so the assessment can be repeated and challenged
		\item[\textbf{(D)}] Record nothing to avoid liability}
	{C}
	{\item Documentation enables review.
		\item Assumptions explain differences over time.
		\item Transparency builds credibility.}
	{Document so it can be reviewed and repeated.}

	% ============================================================
	\section{Risk Response and Reporting (32 percent of the exam)}
	% ============================================================

	\hq{47}{Control Deficiency Severity}
	{Testing finds a control failure. How should its severity be determined?}
	{\item[\textbf{(A)}] By considering the likelihood and magnitude of the potential impact and the presence of compensating controls
		\item[\textbf{(B)}] By the length of the test report
		\item[\textbf{(C)}] By who found it
		\item[\textbf{(D)}] By the tester's mood}
	{A}
	{\item Severity reflects risk, not volume.
		\item Compensating controls can lower impact.
		\item Consistent criteria improve prioritization.}
	{Severity = risk-based judgment.}

	\hq{48}{Amber and Red Thresholds}
	{A KRI uses green, amber, and red zones. What should happen when it enters amber?}
	{\item[\textbf{(A)}] Ignore it until red
		\item[\textbf{(B)}] Investigate the trend and prepare responses before it reaches red, per the defined procedure
		\item[\textbf{(C)}] Hide the indicator
		\item[\textbf{(D)}] Reset the zones}
	{B}
	{\item Amber is an early warning.
		\item Early action prevents breaches.
		\item Procedures define who acts.}
	{Amber is time to act, not time to wait.}

	\hq{49}{Metrics Gaming}
	{Teams close tickets quickly by marking them resolved without fixing the cause. What is the \textbf{BEST} response?}
	{\item[\textbf{(A)}] Add more ticket targets
		\item[\textbf{(B)}] Reward speed only
		\item[\textbf{(C)}] Ignore it
		\item[\textbf{(D)}] Review metric design to reward real outcomes, and add quality checks such as reopen rates}
	{D}
	{\item What gets measured shapes behavior.
		\item Quality checks counter gaming.
		\item Outcome-based measures are more meaningful.}
	{Measure outcomes to avoid gamed metrics.}

	\hq{50}{Audit Committee Reporting}
	{What should risk reporting to the audit committee emphasize?}
	{\item[\textbf{(A)}] Technical logs
		\item[\textbf{(B)}] Significant risks, control weaknesses, remediation progress, and independent assurance results
		\item[\textbf{(C)}] Minor incidents only
		\item[\textbf{(D)}] Staff birthdays}
	{B}
	{\item The committee oversees risk and controls.
		\item Progress shows accountability.
		\item Material issues need attention.}
	{Give the committee what it needs to oversee.}

	\hq{51}{Crisis Communication}
	{During a major incident, which communication practice is \textbf{BEST}?}
	{\item[\textbf{(A)}] Provide timely, accurate, coordinated messages through designated spokespeople
		\item[\textbf{(B)}] Let everyone speak freely
		\item[\textbf{(C)}] Say nothing
		\item[\textbf{(D)}] Provide speculation}
	{A}
	{\item Consistent messages protect trust.
		\item Designated spokespeople avoid mixed statements.
		\item Accuracy matters more than speed alone.}
	{One voice, accurate facts.}

	\hq{52}{Vendor Scorecards}
	{What is the benefit of vendor performance and risk scorecards?}
	{\item[\textbf{(A)}] They replace contracts
		\item[\textbf{(B)}] They hide issues
		\item[\textbf{(C)}] They provide consistent, ongoing visibility into vendor performance against service and control expectations
		\item[\textbf{(D)}] They end oversight}
	{C}
	{\item Consistent measures allow comparison.
		\item Trends show deteriorating vendors.
		\item Results support renewal decisions.}
	{Scorecards turn vendor oversight into data.}

	\hq{53}{Contract Renewal}
	{A critical vendor contract is up for renewal. What risk activity should occur?}
	{\item[\textbf{(A)}] Auto-renew without review
		\item[\textbf{(B)}] Ignore past issues
		\item[\textbf{(C)}] Skip the risk review
		\item[\textbf{(D)}] Reassess the vendor's risk, performance, and incidents, and update terms accordingly}
	{D}
	{\item Conditions may have changed.
		\item Renewal is leverage for improvement.
		\item Update security and exit terms.}
	{Renewal is a risk checkpoint.}

	\hq{54}{Testing the Exit Plan}
	{A cloud exit plan exists but has never been exercised. What is the \textbf{BEST} recommendation?}
	{\item[\textbf{(A)}] Assume it works
		\item[\textbf{(B)}] Test the exit plan, at least in part, to confirm that data and services can be moved within acceptable time
		\item[\textbf{(C)}] Delete the plan
		\item[\textbf{(D)}] Never test}
	{B}
	{\item Untested plans may fail.
		\item Testing reveals dependencies.
		\item Results inform contract terms.}
	{Test the exit before you need it.}

	\hq{55}{Change Failure Rate}
	{What does a high change failure rate indicate from a risk perspective?}
	{\item[\textbf{(A)}] Excellent controls
		\item[\textbf{(B)}] No risk
		\item[\textbf{(C)}] Faster innovation only
		\item[\textbf{(D)}] Weak change testing or approval controls, increasing the risk of outages and incidents}
	{D}
	{\item Failures cause disruption.
		\item Process weaknesses may be the cause.
		\item Improvements reduce risk.}
	{Failed changes signal weak change control.}

	\hq{56}{Patch SLA}
	{The organization defines time limits for patching by severity. What is the main benefit?}
	{\item[\textbf{(A)}] It removes the need to patch
		\item[\textbf{(B)}] It sets clear expectations and enables measurement and escalation when deadlines are missed
		\item[\textbf{(C)}] It delays all patches
		\item[\textbf{(D)}] It hides vulnerabilities}
	{B}
	{\item Service levels create accountability.
		\item Measurement shows compliance.
		\item Escalation addresses delays.}
	{Define deadlines, then measure them.}

	\hq{57}{Vulnerability Aging}
	{A report shows many critical vulnerabilities open for over 180 days. What is the \textbf{BEST} interpretation?}
	{\item[\textbf{(A)}] The environment is secure
		\item[\textbf{(B)}] Vulnerabilities fix themselves
		\item[\textbf{(C)}] Remediation processes are not effective, so exposure and likely audit findings are rising
		\item[\textbf{(D)}] Age does not matter}
	{C}
	{\item Aging shows remediation weakness.
		\item Attackers favor known vulnerabilities.
		\item Escalation and resourcing may be needed.}
	{Aged criticals are an alarm bell.}

	\hq{58}{Mean Time to Remediate}
	{Which metric \textbf{BEST} reflects the speed of fixing vulnerabilities?}
	{\item[\textbf{(A)}] Number of scans run
		\item[\textbf{(B)}] Mean time to remediate by severity
		\item[\textbf{(C)}] Number of tools bought
		\item[\textbf{(D)}] Number of dashboards}
	{B}
	{\item It measures outcome speed.
		\item Segmenting by severity prioritizes risk.
		\item Trends show improvement.}
	{Measure how quickly you fix, not how often you look.}

	\hq{59}{Business Case for Controls}
	{How should the risk practitioner \textbf{BEST} support a business case for a new control?}
	{\item[\textbf{(A)}] Use fear and uncertainty only
		\item[\textbf{(B)}] Show vendor slides only
		\item[\textbf{(C)}] Show the risk reduction, costs, alternatives, and alignment with appetite and objectives
		\item[\textbf{(D)}] Skip cost analysis}
	{C}
	{\item Decision makers need balanced information.
		\item Alternatives demonstrate rigor.
		\item Alignment with objectives gains support.}
	{Support controls with evidence and options.}

	\hq{60}{Return on Security Investment}
	{How can a security investment's value be \textbf{BEST} expressed to executives?}
	{\item[\textbf{(A)}] By the number of features
		\item[\textbf{(B)}] As reduction in expected loss and risk exposure relative to cost, along with strategic benefits
		\item[\textbf{(C)}] By the vendor's reputation only
		\item[\textbf{(D)}] By the length of the report}
	{B}
	{\item Executives think in value and exposure.
		\item Expected loss reduction is comparable to cost.
		\item Non-financial benefits add context.}
	{Speak in risk reduced per dollar spent.}

	\hq{61}{Budget Prioritization by Risk}
	{Competing proposals request budget. How should a risk-based allocation work?}
	{\item[\textbf{(A)}] Fund the loudest sponsor
		\item[\textbf{(B)}] Rank proposals by the risk they reduce relative to appetite and cost, then allocate to highest value
		\item[\textbf{(C)}] Split evenly always
		\item[\textbf{(D)}] Fund the newest technology}
	{B}
	{\item Risk reduction guides value.
		\item Appetite sets urgency.
		\item Transparency builds confidence.}
	{Allocate by risk reduction.}

	\hq{62}{Risk Reduction Verification}
	{After a big risk reduction project, how can the organization verify the result?}
	{\item[\textbf{(A)}] Assume success
		\item[\textbf{(B)}] Rely on the project team's statement
		\item[\textbf{(C)}] Skip verification
		\item[\textbf{(D)}] Re-assess the risk using updated evidence, testing, and indicators and compare to the baseline}
	{D}
	{\item Evidence confirms outcomes.
		\item Baseline comparison shows change.
		\item Update the register.}
	{Verify the benefit with evidence.}

	\hq{63}{Testing Independence}
	{Who should \textbf{ideally} test the effectiveness of a key control?}
	{\item[\textbf{(A)}] The control operator alone
		\item[\textbf{(B)}] The vendor selling it
		\item[\textbf{(C)}] Someone independent of the control's operation, such as the second or third line
		\item[\textbf{(D)}] Nobody}
	{C}
	{\item Independence improves objectivity.
		\item Operators may overlook problems.
		\item Self-testing is useful but limited.}
	{Independent testing is more credible.}

	\hq{64}{Sampling Bias}
	{A tester picks only the cleanest-looking records as a sample. What is the issue?}
	{\item[\textbf{(A)}] Selection bias may overstate control effectiveness, so sampling should be representative or random
		\item[\textbf{(B)}] No issue
		\item[\textbf{(C)}] The sample is too accurate
		\item[\textbf{(D)}] It improves quality}
	{A}
	{\item Biased samples mislead.
		\item Random or risk-based sampling is better.
		\item Document the method.}
	{Do not cherry-pick your sample.}

	\hq{65}{SLA Monitoring}
	{A provider's reports show all service levels met, but users report frequent outages. What should the practitioner do?}
	{\item[\textbf{(A)}] Trust the provider fully
		\item[\textbf{(B)}] Ignore users
		\item[\textbf{(C)}] Cancel immediately without facts
		\item[\textbf{(D)}] Validate the provider's measurement methods and compare with independent or user data}
	{D}
	{\item Reports may exclude relevant outages.
		\item Independent data gives perspective.
		\item Contract definitions matter.}
	{Verify what vendors report.}

	\hq{66}{Automated Reporting}
	{What is a risk of heavy automation in risk reporting?}
	{\item[\textbf{(A)}] Reports become too accurate
		\item[\textbf{(B)}] It removes all errors
		\item[\textbf{(C)}] Errors in data or logic can propagate unnoticed, so validation and ownership remain essential
		\item[\textbf{(D)}] It needs no oversight}
	{C}
	{\item Automation scales mistakes.
		\item Data lineage should be understood.
		\item Humans must review exceptions.}
	{Automate the work, not the accountability.}

	\hq{67}{Drill-Down Reports}
	{Why provide drill-down capability in risk dashboards?}
	{\item[\textbf{(A)}] It lets decision makers move from summary to supporting detail for understanding and challenge
		\item[\textbf{(B)}] To confuse the audience
		\item[\textbf{(C)}] To hide data
		\item[\textbf{(D)}] To increase meetings}
	{A}
	{\item Summaries need credible detail.
		\item Detail enables verification.
		\item Different users need different depth.}
	{Summary on top, evidence underneath.}

	\hq{68}{Risk-Adjusted Decisions}
	{Which approach \textbf{BEST} reflects risk-adjusted decision making for a new initiative?}
	{\item[\textbf{(A)}] Consider only benefits
		\item[\textbf{(B)}] Consider only risks
		\item[\textbf{(C)}] Evaluate expected benefits against the associated risks, including residual risk, before approval
		\item[\textbf{(D)}] Decide by seniority}
	{C}
	{\item Balanced evaluation leads to better outcomes.
		\item Residual risk should be known.
		\item Owners approve within appetite.}
	{Weigh benefits against risk, not instead of it.}

	\hq{69}{Third-Party Breach Response}
	{A vendor breach exposes your employee data. What should the risk practitioner prioritize?}
	{\item[\textbf{(A)}] Wait for the news
		\item[\textbf{(B)}] Rate the vendor as perfect
		\item[\textbf{(C)}] Assess impact, ensure containment and notification obligations are met, and update the vendor risk
		\item[\textbf{(D)}] Do nothing}
	{C}
	{\item Obligations may be time-sensitive.
		\item Impact informs response.
		\item Findings update vendor tiering.}
	{Respond to the incident, then learn from it.}

	\hq{70}{Control Framework Mapping}
	{The organization must comply with several standards. What is the benefit of mapping controls across frameworks?}
	{\item[\textbf{(A)}] It reduces duplication by showing where one control satisfies multiple requirements
		\item[\textbf{(B)}] It creates more controls
		\item[\textbf{(C)}] It removes compliance duties
		\item[\textbf{(D)}] It hides gaps}
	{A}
	{\item Efficiency improves.
		\item Common controls simplify testing.
		\item Gaps are easier to see.}
	{Map once, comply many times.}

	\hq{71}{Testing Frequency}
	{How should the frequency of control testing be determined?}
	{\item[\textbf{(A)}] Same for all controls
		\item[\textbf{(B)}] Only when convenient
		\item[\textbf{(C)}] Based on the control's criticality, risk level, and previous results
		\item[\textbf{(D)}] Only after incidents}
	{C}
	{\item Higher risk needs more frequent testing.
		\item Previous failures justify closer attention.
		\item Resources are focused.}
	{Test the important controls more often.}

	\hq{72}{Exception Register}
	{Why maintain an exception register with expiry dates?}
	{\item[\textbf{(A)}] To ensure exceptions are visible, reviewed, and not allowed to become permanent unmanaged risk
		\item[\textbf{(B)}] To prove everyone follows policy
		\item[\textbf{(C)}] To delay all reviews
		\item[\textbf{(D)}] To remove owners}
	{A}
	{\item Visibility supports oversight.
		\item Expiry forces review.
		\item Compensating controls can be verified.}
	{Exceptions should expire.}

	\hq{73}{Escalating Under Pressure}
	{The business pressures the practitioner not to escalate a major risk before a product launch. What should be done?}
	{\item[\textbf{(A)}] Escalate according to policy, presenting facts and options so leadership can decide knowingly
		\item[\textbf{(B)}] Stay silent
		\item[\textbf{(C)}] Delete the risk
		\item[\textbf{(D)}] Delay until after launch}
	{A}
	{\item Material risk requires timely escalation.
		\item Facts support decisions.
		\item Leadership can accept risk deliberately.}
	{Do not let launch dates hide risks.}

	\hq{74}{Single Version of Truth}
	{Different reports give different risk numbers for the same area. What is the \textbf{BEST} fix?}
	{\item[\textbf{(A)}] Allow every report to differ
		\item[\textbf{(B)}] Establish authoritative data sources, definitions, and a single reporting process
		\item[\textbf{(C)}] Hide reports
		\item[\textbf{(D)}] Pick numbers randomly}
	{B}
	{\item Consistent data supports trust.
		\item Definitions avoid disputes.
		\item Single processes reduce errors.}
	{One set of numbers.}

	\hq{75}{Post-Incident Report}
	{What should a post-incident report to management primarily contain?}
	{\item[\textbf{(A)}] Blame assignments only
		\item[\textbf{(B)}] Technical commands only
		\item[\textbf{(C)}] A single sentence
		\item[\textbf{(D)}] Root cause, impact, response effectiveness, lessons learned, and corrective actions with owners}
	{D}
	{\item Reports drive improvement.
		\item Owners and dates ensure follow-up.
		\item Lessons update the risk register.}
	{Close the loop with actions and owners.}

	\hq{76}{Business Continuity Metrics}
	{Which metric \textbf{BEST} indicates continuity readiness?}
	{\item[\textbf{(A)}] Number of plans printed
		\item[\textbf{(B)}] Number of staff trained once
		\item[\textbf{(C)}] Results of exercises against recovery objectives, and the closure of identified gaps
		\item[\textbf{(D)}] Plan page count}
	{C}
	{\item Exercises prove capability.
		\item Gap closure shows improvement.
		\item Objectives tie to business needs.}
	{Readiness is demonstrated, not documented.}

	\hq{77}{Policy Violation Trends}
	{Reports show increasing policy violations in one department. What should the practitioner do?}
	{\item[\textbf{(A)}] Ignore it
		\item[\textbf{(B)}] Analyze root causes with the owner and address training, process, or incentive issues
		\item[\textbf{(C)}] Fire the department
		\item[\textbf{(D)}] Delete the data}
	{B}
	{\item Trends signal systemic problems.
		\item Root causes may be non-obvious.
		\item Targeted fixes work better than blame.}
	{Trends point to systemic causes.}

	\hq{78}{Closing Risk Issues}
	{When may a risk issue be closed in the register?}
	{\item[\textbf{(A)}] When the owner stops responding
		\item[\textbf{(B)}] When response actions are complete and verified, and the residual risk is within appetite or formally accepted
		\item[\textbf{(C)}] When the deadline passes
		\item[\textbf{(D)}] When the report is filed}
	{B}
	{\item Closure needs evidence.
		\item Residual risk must be within appetite or accepted.
		\item Verification prevents false closure.}
	{Close risks on evidence, not on dates.}

	% ============================================================
	\section{Information Technology and Security (22 percent of the exam)}
	% ============================================================

	\hq{79}{Identity Lifecycle}
	{Which practice \textbf{BEST} manages user access throughout an employee's tenure?}
	{\item[\textbf{(A)}] Grant access once and never review
		\item[\textbf{(B)}] Joiner, mover, and leaver processes tied to HR events, with periodic access reviews
		\item[\textbf{(C)}] Let users request any access by phone
		\item[\textbf{(D)}] Share accounts in teams}
	{B}
	{\item Job changes alter access needs.
		\item HR events are reliable triggers.
		\item Reviews catch accumulated privileges.}
	{Access should follow the person's role over time.}

	\hq{80}{Privileged Access Management}
	{Which capability is \textbf{MOST} valuable in privileged access management?}
	{\item[\textbf{(A)}] Permanent shared admin passwords
		\item[\textbf{(B)}] Unlimited sessions without logs
		\item[\textbf{(C)}] Just-in-time elevation, credential vaulting, and session recording
		\item[\textbf{(D)}] Local admin for everyone}
	{C}
	{\item Limiting standing privilege reduces exposure.
		\item Vaulting protects credentials.
		\item Recording supports accountability.}
	{Make privilege temporary, controlled, and recorded.}

	\hq{81}{Session Management}
	{Users leave sensitive applications logged in on shared computers. Which control \textbf{BEST} reduces risk?}
	{\item[\textbf{(A)}] Automatic session timeouts and re-authentication for sensitive actions
		\item[\textbf{(B)}] Longer sessions
		\item[\textbf{(C)}] No timeouts
		\item[\textbf{(D)}] Posters about locking screens only}
	{A}
	{\item Timeouts limit unauthorized use.
		\item Step-up authentication protects critical actions.
		\item Technical controls are more reliable than reminders.}
	{Do not rely on users to log out.}

	\hq{82}{Network Access Control}
	{Unknown personal devices connect to the corporate network. What \textbf{BEST} addresses this?}
	{\item[\textbf{(A)}] Disable all logging
		\item[\textbf{(B)}] Print the Wi-Fi password on walls
		\item[\textbf{(C)}] Allow all devices
		\item[\textbf{(D)}] Network access control that authenticates devices and checks their compliance before granting access}
	{D}
	{\item Device identity and posture matter.
		\item Non-compliant devices can be isolated.
		\item Visibility supports inventory.}
	{Authenticate the device, not only the user.}

	\hq{83}{Microsegmentation}
	{After a breach, attackers moved freely among servers. Which approach \textbf{BEST} limits future lateral movement?}
	{\item[\textbf{(A)}] A flat network
		\item[\textbf{(B)}] More shared credentials
		\item[\textbf{(C)}] Microsegmentation with least-privilege network policies between workloads
		\item[\textbf{(D)}] Fewer logs}
	{C}
	{\item Fine-grained segments contain compromise.
		\item Policies restrict necessary paths only.
		\item Monitoring detects unusual flows.}
	{Contain the blast radius.}

	\hq{84}{Security Automation}
	{The security team is overwhelmed by repetitive tasks. Which capability \textbf{BEST} helps?}
	{\item[\textbf{(A)}] Security orchestration, automation, and response for routine, well-defined workflows
		\item[\textbf{(B)}] Hiring no one
		\item[\textbf{(C)}] Ignoring alerts
		\item[\textbf{(D)}] Disabling detection}
	{A}
	{\item Automation handles repeat tasks.
		\item Analysts focus on complex work.
		\item Playbooks must be reviewed.}
	{Automate the routine; keep humans for judgment.}

	\hq{85}{Threat Hunting}
	{What is the primary goal of proactive threat hunting?}
	{\item[\textbf{(A)}] To wait for alerts
		\item[\textbf{(B)}] To replace all tools
		\item[\textbf{(C)}] To reduce logging
		\item[\textbf{(D)}] To search for hidden threats that automated detection may have missed}
	{D}
	{\item Hunting assumes compromise is possible.
		\item Hypotheses guide searches.
		\item Findings improve detections.}
	{Hunt for what the alerts did not catch.}

	\hq{86}{Operational Technology}
	{A plant's control systems cannot be patched frequently without risking production. What is the \textbf{BEST} approach?}
	{\item[\textbf{(A)}] Treat it exactly like office IT
		\item[\textbf{(B)}] Ignore it
		\item[\textbf{(C)}] Use network segmentation, access control, monitoring, and compensating controls tailored to safety and availability
		\item[\textbf{(D)}] Connect it directly to the internet}
	{C}
	{\item OT prioritizes safety and availability.
		\item Segmentation reduces exposure.
		\item Monitoring detects anomalies.}
	{OT security protects safety first.}

	\hq{87}{Mobile App Security}
	{A company builds a mobile app handling sensitive data. Which practice is \textbf{MOST} important?}
	{\item[\textbf{(A)}] Secure development, secure data storage, strong authentication, and regular testing
		\item[\textbf{(B)}] A colorful icon
		\item[\textbf{(C)}] Larger downloads
		\item[\textbf{(D)}] Fewer updates}
	{A}
	{\item Mobile devices can be lost or compromised.
		\item Data storage should be minimized and protected.
		\item Testing finds weaknesses.}
	{Build security into the app lifecycle.}

	\hq{88}{API Protection}
	{Public APIs are being abused with excessive requests. Which control \textbf{BEST} helps?}
	{\item[\textbf{(A)}] Disable authentication
		\item[\textbf{(B)}] Authentication, authorization, rate limiting, and monitoring through an API gateway
		\item[\textbf{(C)}] Unlimited requests
		\item[\textbf{(D)}] Public documentation only}
	{B}
	{\item Gateways enforce policies.
		\item Rate limits reduce abuse.
		\item Monitoring shows misuse.}
	{Control who calls the API and how often.}

	\hq{89}{Authorization Standards}
	{An application lets users sign in with another provider and grants limited access to their data. Which standard commonly supports delegated authorization?}
	{\item[\textbf{(A)}] OAuth
		\item[\textbf{(B)}] A cabling standard
		\item[\textbf{(C)}] A fax protocol
		\item[\textbf{(D)}] A paper format}
	{A}
	{\item OAuth delegates limited access.
		\item Scopes limit what is granted.
		\item Tokens should be protected.}
	{OAuth delegates access without sharing passwords.}

	\hq{90}{SQL Injection}
	{Which practice \textbf{BEST} prevents SQL injection?}
	{\item[\textbf{(A)}] Longer passwords
		\item[\textbf{(B)}] Bigger databases
		\item[\textbf{(C)}] Parameterized queries with input validation
		\item[\textbf{(D)}] More servers}
	{C}
	{\item Parameters separate data from commands.
		\item Validation reduces risk further.
		\item Least privilege limits impact.}
	{Never build queries by concatenating user input.}

	\hq{91}{Code Signing}
	{What is the main purpose of code signing?}
	{\item[\textbf{(A)}] To verify the publisher and that the code has not been altered since signing
		\item[\textbf{(B)}] To make code run faster
		\item[\textbf{(C)}] To encrypt all data
		\item[\textbf{(D)}] To reduce file size}
	{A}
	{\item Signatures provide authenticity and integrity.
		\item Private keys must be protected.
		\item Verification happens at installation or run time.}
	{Signed code proves origin and integrity.}

	\hq{92}{Software Bill of Materials}
	{Why is a software bill of materials valuable?}
	{\item[\textbf{(A)}] It replaces testing
		\item[\textbf{(B)}] It lists components in software so that vulnerable ones can be identified quickly when issues are disclosed
		\item[\textbf{(C)}] It guarantees security
		\item[\textbf{(D)}] It lists employees}
	{B}
	{\item Visibility enables fast response.
		\item Supplier components are included.
		\item It supports supply chain risk management.}
	{Know the ingredients of your software.}

	\hq{93}{Key Rotation}
	{Why rotate encryption keys periodically?}
	{\item[\textbf{(A)}] To make keys longer
		\item[\textbf{(B)}] To avoid backups
		\item[\textbf{(C)}] To confuse administrators
		\item[\textbf{(D)}] To limit the amount of data exposed by a compromised key and to meet policy or regulatory requirements}
	{D}
	{\item Rotation limits exposure windows.
		\item Compromise may be undetected.
		\item Procedures must handle old data.}
	{Rotate keys to limit exposure.}

	\hq{94}{Data Masking}
	{Analysts need realistic data for reports but should not see actual personal details. Which control fits?}
	{\item[\textbf{(A)}] Share raw data widely
		\item[\textbf{(B)}] Data masking or tokenization of sensitive fields
		\item[\textbf{(C)}] Email full databases
		\item[\textbf{(D)}] Remove all access controls}
	{B}
	{\item Masking reduces exposure.
		\item Tokenization replaces values with surrogates.
		\item Access to originals stays restricted.}
	{Mask what analysts do not need to see.}

	\hq{95}{Secure File Transfer}
	{Sensitive files are regularly emailed to partners. What is the \textbf{BEST} improvement?}
	{\item[\textbf{(A)}] Continue emailing
		\item[\textbf{(B)}] Use public file sharing without controls
		\item[\textbf{(C)}] Use secure, encrypted file transfer with authentication and logging
		\item[\textbf{(D)}] Print and mail}
	{C}
	{\item Email is often unencrypted.
		\item Controlled transfers provide traceability.
		\item Authentication ensures only intended recipients.}
	{Use controlled channels for sensitive data.}

	\hq{96}{Software License Compliance}
	{Audits show the company uses more software licenses than it has purchased. What risk does this create?}
	{\item[\textbf{(A)}] No risk
		\item[\textbf{(B)}] Better security
		\item[\textbf{(C)}] Faster performance
		\item[\textbf{(D)}] Legal and financial exposure from non-compliance with license terms}
	{D}
	{\item Vendors can audit and claim penalties.
		\item An inventory supports compliance.
		\item Governance should track entitlements.}
	{Unlicensed use creates legal and financial risk.}

	\hq{97}{Virtualization Sprawl}
	{Virtual machines are created without lifecycle management and many are forgotten. What is the risk?}
	{\item[\textbf{(A)}] Lower attack surface
		\item[\textbf{(B)}] Better monitoring
		\item[\textbf{(C)}] No impact
		\item[\textbf{(D)}] Unpatched, unmonitored systems that increase attack surface and cost}
	{D}
	{\item Forgotten systems escape controls.
		\item Inventory and lifecycle processes help.
		\item Decommissioning reduces exposure.}
	{Forgotten VMs are unmanaged risk.}

	\hq{98}{Hypervisor Security}
	{Why is hypervisor security especially important?}
	{\item[\textbf{(A)}] It affects only one VM
		\item[\textbf{(B)}] It has no connection to guests
		\item[\textbf{(C)}] A compromised hypervisor could affect all the virtual machines it hosts
		\item[\textbf{(D)}] It only affects printing}
	{C}
	{\item The hypervisor is a high-value target.
		\item Hardening and patching are critical.
		\item Access should be tightly restricted.}
	{Protect the layer beneath all the VMs.}

	\hq{99}{Cloud Exit Data}
	{When ending a cloud service, what should be verified regarding data?}
	{\item[\textbf{(A)}] That all data has been returned or migrated and then securely deleted by the provider, with confirmation
		\item[\textbf{(B)}] That the provider keeps the data indefinitely
		\item[\textbf{(C)}] That nothing happens
		\item[\textbf{(D)}] That logs are disabled}
	{A}
	{\item Data may remain in provider systems.
		\item Contracts should specify deletion.
		\item Evidence of deletion supports compliance.}
	{Verify deletion when you leave.}

	\hq{100}{Tokenization of Payment Data}
	{A merchant wants to reduce the scope and risk of storing card numbers. Which control \textbf{BEST} helps?}
	{\item[\textbf{(A)}] Storing card numbers in a spreadsheet
		\item[\textbf{(B)}] Emailing card numbers to staff
		\item[\textbf{(C)}] Writing numbers on tickets
		\item[\textbf{(D)}] Tokenization, replacing card numbers with non-sensitive tokens held by a secure provider}
	{D}
	{\item Tokens have no exploitable value on their own.
		\item Less stored card data means a smaller compliance scope.
		\item The token vault must be strongly protected.}
	{If you do not store it, you cannot lose it.}

	% ============================================================
	\section{Conclusion}
	% ============================================================

	\begin{frame}{Exam Mindset Summary}
		\begin{block}{Remember These Core Principles}
			\begin{itemize}
				\item \textbf{Business First:} Risk is defined by its effect on business objectives.
				\item \textbf{Owners Decide:} Practitioners inform; owners accept, mitigate, transfer, or avoid.
				\item \textbf{Three Lines:} The first line owns risk, the second oversees, and the third gives independent assurance.
				\item \textbf{Ethics Matter:} Honesty, independence, and confidentiality underpin credible risk advice.
				\item \textbf{Evidence Over Opinion:} Use data, scenarios, and calibrated ratings.
				\item \textbf{Proportionality:} Controls should fit the risk and its cost.
				\item \textbf{Accountability Cannot be Outsourced:} Vendors and cloud providers share duties, but accountability remains with the organization.
			\end{itemize}
		\end{block}
	\end{frame}

	\begin{frame}{Thank You}
		\begin{center}
			{\Huge \textbf{Good Luck on Your CRISC Exam!}}

			\vspace{1cm}

			\normalsize
			Real exam questions test your ability to apply these concepts in practical, real-world scenarios.
		\end{center}
	\end{frame}

\end{document}
